% =====================================================================
% EJE A - EL PROBLEMA / DOMINIO  ---  RESPONSABLE: INTEGRANTE 1
% Como ha abordado la literatura la prediccion de resultados en
% carreras y deportes con datos tabulares? Que enfoques dominan?
% Que se considera resuelto y que sigue abierto?
% ORGANIZAR POR ENFOQUE, no un parrafo por articulo.
% =====================================================================
\subsection{Axis A: The Problem Domain}
\label{sec:ejeA}

Predicting the outcome of a horse race is framed in the literature in two largely
disconnected ways: as a supervised classification problem over pre-race tabular
features, and, less frequently, as a market-efficiency question about whether
bookmaker odds already encode most of the predictable signal. These objectives are
not interchangeable: they involve different target variables, evaluation procedures,
and interpretations of predictive performance. Our review found supervised learning
over pre-race tabular information is a recurring approach, while the formulation of
the race as a direct within-race ranking problem is less explicitly addressed.

\subsubsection*{\textbf{Classical statistical models}}
Early and still-recurring approaches rely on
regression rather than modern ML. Li et al.\cite{ejeA_li2025} fit linear and
log-linear regression models on going (track condition), distance and class to
predict winning times at three UK racecourses, reporting $R^2 \approx 0.995$ for
the fitted times themselves. However, the objective of the study is to audit
whether officially reported track conditions are consistent with observed
performance, not to predict the winner or establish a ranking of the horses; they
explicitly note this is a starting point and suggest tree-based ML as a natural
extension. This distinction isimportant because it illustrates a recurring pattern:
regression-based work in this domain has mostly targeted auxiliary quantities
(going, finishing time) rather than the win/place outcome itself, which is the kind
of outcome-prediction aimed to solved in this project.

\subsubsection*{\textbf{Machine learning approaches}} The reviewed horse-racing studies show a
range of supervised learning approaches and experimental settings. Davoodi and
Khanteymoori~\cite{ejeA_davoodi2010} use artificial neural networks to predict the
finishing time of individual horses at Aqueduct Race Track in New York. Their
dataset contains 100 races from January 2010, and the authors compare five
supervised training algorithms. They report an average accuracy of approximately
77\%. The study therefore provides evidence that nonlinear supervised models can be
applied to horse-racing data, but its small temporal and geographic scope limits
conclusions about generalization.

Schumaker\cite{ejeA_schumaker2018} studies harness racing at Northfield Park in
Ohio using Support Vector Regression and historical race-performance variables. A
key contribution of the study is its explicit temporal separation: data from
October 2009 through September 2010 are used for training, while October through
December 2010 are reserved for testing. The study also compares different race
history windows and reports that the Time 4--3 configuration performs substantially
better than the Time 7--3 configuration for the harness-racing setting. In
particular, Time 7--3 achieves 35.42\% accuracy, whereas Time 4--3 achieves 50.44\%
across the reported wagering analysis. These results shows that feature-history
configurations developed for one racing subtype should not be assumed to transfer
directly to another without validation. An important clarifition is that with the
term Time 4--3, for example, we refer to the previous races being considered as
reference for the study settings, in this case 4--3 considers from 4th to 3th
previous race.

Gupta and Singh~\cite{ejeA_gupta2023} examine horse-race prediction using data
collected from a horse-racing website in India. Their dataset contains 14,750
records and 56 attributes, with wins representing only 11.8\% of the observations.
The authors report that a Random Forest model trained without imbalance handling
can achieve high accuracy while failing to identify the minority class, whereas
applying SMOTE improves the reported classification results to 93.1\% accuracy,
92.9\% F1-score (balance between precision and recall), and 97.6\% ROC-AUC (ability
to distinguish between classes across different decision thresholds). This study
highlights an important domain characteristic: accuracy alone can be misleading
when the target distribution is highly imbalanced, making class-sensitive metrics
and appropriate training procedures important for horse-racing prediction. A
related study by Gupta and Singh~\cite{ejeA_gupta2024b} compares K-nearest
neighbors (KNN), logistic regression, Random Forest, Gaussian Naive Bayes,
AdaBoost, and Bagging, together with Information Gain and Chi-Square feature-
selection procedures. The authors report that the effectiveness of feature
selection depends on the selection strategy, although, the Random Forest provides
the best performance among the evaluated models across the despite feature-
selection settings, followed by AdaBoost as another good option.

\subsubsection*{\textbf{Market efficiency and odds as a predictor}} Colle~\cite{ejeA_colle_nd}
notes that horse-racing prediction research is comparatively scarce next to
mainstream team sports precisely because commercial betting syndicates keep their
methods private, and traces the field's statistical roots to pari-mutuel market
modelling from the 1980s. Chavan et al.~\cite{ejeC_chavan2024} summarize a study on
Champ de Mars (Mauritius) in which betting odds alone reached only around 7.4\%
correct-rank prediction, which they interpret as evidence that odds are not, by
themselves, a strong predictor of finishing order on that track.

\subsubsection*{\textbf{What is comparatively settled}} Across the studies we read, tree-ensemble
and kernel methods consistently outperform plain linear/logistic baselines on
tabular horse-racing features once class imbalance is addressed, and feature sets
built from pre-race information (weight, draw, jockey/trainer history, ratings)
are treated as the standard input space; nobody in our corpus argues that these
features are irrelevant. What remains open is precisely what our project
targets: first, every horse-racing-specific study we found evaluates on a single
track, a short time window (weeks to a few years) or a single jurisdiction, so
generalization across time and venues is essentially untested; second, none of the
studies we read formulate the problem as a within-race ranking task (predicting an
order over the runners of a given race) rather than an independent per-horse
binary classification, even though a race is intrinsically a comparison among a
fixed, known group of competitors; and third, explicit prevention of temporal data
leakage is rare. Only Schumaker~\cite{ejeA_schumaker2018} uses a clean
chronological split, while the others do not describe how (or whether) they
avoided mixing future information into training folds. These three gaps are
picked up again, together with the dataset-specific ones, in
Section~\ref{sec:brechas}.

\subsubsection*{\textbf{Synthesis of the domain}} The reviewed literature does not
support the claim that a single modeling approach has been established as universally
superior for horse-racing prediction. Instead, the studies demonstrate that the problem
has been formulated in different ways, including continuous prediction of performance
variables, individual-horse outcome prediction, and wagering-oriented prediction.
Tree-based ensembles, neural networks, support-vector methods, and linear models
all appear in the reviewed literature, but their reported results depend strongly
on the target variable, dataset, racing subtype, feature construction, and
evaluation protocol.

Two methodological patterns are particularly relevant to the present project.
First, several studies operate within a single racecourse, jurisdiction, or
relatively short observation period, which limits the evidence available about
generalization across time and venues. Second, the reviewed studies use different
target definitions and evaluation procedures, making direct comparisons of reported
accuracy difficult. Schumaker~\cite{ejeA_schumaker2018} provides a useful
counterexample by explicitly separating historical training data from a later
testing period, illustrating the importance of evaluating a model under a temporal
setting when the objective is future prediction.

These observations motivate two research gaps for the present project. First, the
generalization of horse-racing prediction models across a longer historical period
and a different data setting remains insufficiently characterized in the reviewed
literature. Second, although race outcomes are inherently comparative because
several horses compete in the same event, the reviewed studies primarily evaluate
individual-horse predictions or wagering decisions rather than explicitly
formulating the task as a within-race ranking problem. These gaps provide the
domain-level motivation for the following axes. The techniques axis evaluates
which classical machine-learning methods are appropriate for structured racing data
and under what assumptions, while the dataset axis examines prior work using the
selected Hong Kong racing data and establishes the most appropriate baseline.
Together, these analyses will determine which aspects of the problem can be
addressed by the proposed study and which remain outside its scope.
